<% const markdownIndent = (latexcompiler == "both") ? '  ' : ''; -%>
% Write text in Markdown instead of LaTeX: \begin{markdown} ... \end{markdown}, or \markdownInput{chapter.md} for a whole file.
% Doc: https://ctan.org/pkg/markdown
% LuaLaTeX only: under pdfLaTeX, the package needs --shell-escape.
<% if (latexcompiler == "both") { -%>
\ifluatex
<% } -%>
<%= markdownIndent %>\usepackage[
<%= markdownIndent %>  citations,          % [@key] citation syntax
<%= markdownIndent %>  hybrid,             % pass raw LaTeX through
<%= markdownIndent %>  notes,              % [^label] referenced footnotes
<%= markdownIndent %>  inlineNotes,        % ^[inline footnote text]
<%= markdownIndent %>  pipeTables,         % | a | b | style tables
<%= markdownIndent %>  tableCaptions,      % : Caption text below a table (typeset above it, see below)
<%= markdownIndent %>  tableAttributes,    % {#id} identifier on a table caption
<%= markdownIndent %>  headerAttributes,   % {#id} identifier on a heading
<%= markdownIndent %>  relativeReferences, % <#id> to reference an id defined via {#id} (typeset by <%- cref %>, see below)
<%= markdownIndent %>  hashEnumerators,    % 1) 2) 3) style enumerators
<%= markdownIndent %>]{markdown}
<% if (isPaper) { -%>
<%= markdownIndent %>% The markdown package typesets "#" as \chapter if the class defines it (llncs does); in a paper, "#" is a section
<%= markdownIndent %>% textidote: ignore begin
<%= markdownIndent %>\renewcommand{\markdownRendererHeadingOne}[1]{\section{#1}}
<%= markdownIndent %>\renewcommand{\markdownRendererHeadingTwo}[1]{\subsection{#1}}
<%= markdownIndent %>\renewcommand{\markdownRendererHeadingThree}[1]{\subsubsection{#1}}
<%= markdownIndent %>% textidote: ignore end
<% } -%>
<%= markdownIndent %>% Table captions above the table, as for all other tables (the markdown package puts them below)
<%= markdownIndent %>\makeatletter
<%= markdownIndent %>\ExplSyntaxOn
<%= markdownIndent %>\markdownSetup{rendererPrototypes={
<%= markdownIndent %>      table = {
<%= markdownIndent %>          \markdownLaTeXTable={}
<%= markdownIndent %>          \markdownLaTeXTableAlignment={}
<%= markdownIndent %>          \markdownLaTeXTableEnd={\markdownLaTeXBottomRule\end{tabular}}
<%= markdownIndent %>          \tl_if_empty:nF {#1} {
<%= markdownIndent %>            \addto@hook\markdownLaTeXTable{\begin{table}\centering\caption{#1}}
<%= markdownIndent %>            \seq_map_inline:Nn \l__markdown_table_identifiers_seq { \addto@hook\markdownLaTeXTable{\label{##1}} }
<%= markdownIndent %>            \addto@hook\markdownLaTeXTableEnd{\end{table}}
<%= markdownIndent %>          }
<%= markdownIndent %>          \addto@hook\markdownLaTeXTable{\begin{tabular}}
<%= markdownIndent %>          \markdownLaTeXRowCounter=0
<%= markdownIndent %>          \markdownLaTeXRowTotal=#2
<%= markdownIndent %>          \markdownLaTeXColumnTotal=#3
<%= markdownIndent %>          \markdownLaTeXRenderTableRow
<%= markdownIndent %>        }
<%= markdownIndent %>    }}
<%= markdownIndent %>\ExplSyntaxOff
<%= markdownIndent %>\makeatother
<%- include('markdown-references.preamble.en.tex', this); -%>
<% if (latexcompiler == "both") { -%>
\fi
<% } -%>
